\documentclass[11pt]{article}

\usepackage{dad}

\begin{document}

\courseheader
  {Data and Design with Python}
  {LMTH 2080}
  {Fall 2026 \textbullet{} Mondays, 9:00--11:40 AM}
  {Eugene Lang College of Liberal Arts}

\begin{tabularx}{\textwidth}{@{}>{\bfseries}l X >{\bfseries}l X@{}}
Instructor & Jacob Koehler &
Email & \href{mailto:koehlerj@newschool.edu}{koehlerj@newschool.edu} \\
Office Hours & By appointment &
Course Meeting & Mondays, 9:00--11:40 AM \\
\end{tabularx}

\section{Course Description}

Data is everywhere, but having data is not the same as having useful data.
This course introduces Python through the problems that data can help us
investigate, explain, and address. We will begin by asking what questions
matter to us, what evidence might help answer them, where that evidence comes
from, and what may be missing from the data we find.

Python will be our primary computational medium. Students will learn to
organize computational projects, write and debug programs, work with tabular
data, acquire data from the web, create visualizations, and build interactive
applications. Later in the semester, we will study basic ideas from machine
learning and neural networks before working with pretrained models from
Hugging Face.

The course culminates in the design and deployment of an interactive
Streamlit application built around a pretrained model. Throughout the
semester, our emphasis will be not only on whether code works, but whether the
data, model, representation, and interface are useful for the problem and
people for whom they are intended.

No prior experience with Python is required. Students are expected to have
access to a laptop computer and bring it to every class.

\begin{keyidea}
The central question of the course is not simply ``What can we do with
Python?'' but \emph{What are we trying to understand or change, and what data
and computational tools would actually help?}
\end{keyidea}

\section{Learning Outcomes}

By the successful completion of this course, students will be able to:

\begin{enumerate}[leftmargin=*]
  \item formulate questions that can be investigated with data and evaluate
        whether a dataset is useful for those questions;
  \item navigate a filesystem from the command line and use Git to record and
        manage changes to a computational project;
  \item write, read, test, and debug Python programs using variables, data
        types, collections, conditionals, loops, and functions;
  \item use \texttt{pandas} to load, inspect, clean, transform, summarize, and
        analyze tabular data;
  \item create and critique data visualizations with attention to visual
        encoding, audience, uncertainty, and purpose;
  \item acquire and structure data from web APIs and other online sources;
  \item build interactive data applications using Streamlit;
  \item explain the basic logic of supervised machine learning, including
        features, targets, training and test data, prediction, evaluation,
        and overfitting;
  \item distinguish model training from inference and explain at a conceptual
        level how pretrained models are used;
  \item evaluate a pretrained model using its documentation, intended uses,
        inputs and outputs, limitations, and observed failure cases; and
  \item design, build, document, and present an interactive application that
        makes thoughtful use of data and a pretrained model.
\end{enumerate}

\section{Course Questions}

We will return to a small set of questions throughout the semester:

\begin{dataprompt}
What problem or question are you interested in? What would you need to know
to investigate it? What data could provide that evidence? Who collected the
data, for what purpose, and what does it fail to capture?
\end{dataprompt}

\begin{designprompt}
Who is this for? What should they be able to see, understand, decide, or do?
How do the choices we make in code, representation, and interface shape that
experience?
\end{designprompt}

\section{Projects and Assignments}

The semester is organized around three increasingly public forms of
computational work:

\subsection{Project I: Data Investigation}

Students will identify or work with a dataset, clean and transform it, ask
questions of it, and communicate findings through analysis and visualization.
The primary artifact will be a well-organized computational notebook.

\subsection{Project II: Designed Data Application}

Students will build a small Streamlit application around a dataset or data
source. The emphasis is on moving from analysis performed for oneself to an
interactive object designed for another person.

\subsection{Final Project: Pretrained Model Application}

Students will select a pretrained model from the Hugging Face Model Hub and
design an interactive Streamlit application around it. Students will be
expected to understand the model well enough to describe its purpose, inputs,
outputs, documentation, known limitations, and observed failure cases. The
final project will be presented in class.

\begin{center}
\large\bfseries
Notebook $\longrightarrow$ Data Application $\longrightarrow$ Model Application
\end{center}

Shorter assignments, readings, coding exercises, and in-class studio work will
support these projects.

\section{Evaluation}

\begin{tabularx}{0.72\textwidth}{@{}X r@{}}
\toprule
\textbf{Component} & \textbf{Weight} \\
\midrule
Active Participation / Attendance & 30\% \\
Homework and Short Assignments & 40\% \\
Projects & 30\% \\
\midrule
\textbf{Total} & \textbf{100\%} \\
\bottomrule
\end{tabularx}

\vspace{6pt}

Projects will be evaluated not only on technical correctness. Depending on
the assignment, evaluation may also consider the quality of the question,
choice and understanding of data, clarity of analysis, design decisions,
documentation, testing, communication, and reflection.

\section{Working With Generative AI}

Generative AI systems can produce code, explanations, and interfaces, but
using generated code effectively requires being able to read it, reason about
it, test it, and recognize when it is wrong.

Early in the semester, some assignments may restrict the use of generative AI
so that everyone develops a working foundation in Python and debugging. As
the course progresses---and particularly as we study machine learning,
transformers, and pretrained models---we will examine these tools more
directly and use them where appropriate.

When AI-assisted tools are permitted, students remain responsible for every
part of the work they submit. You should be able to explain your code, modify
it, test it, identify its assumptions, and describe how AI tools contributed
to the work. Assignment-specific instructions take precedence over this
general policy.

\section{Materials and Computing}

Students should bring a laptop to every class. We will use a Unix-like command
line, Git and GitHub, Python, Jupyter, \texttt{pandas}, visualization
libraries, Streamlit, scikit-learn, and Hugging Face tools.

Software used in the course is freely available. Installation instructions
and course-specific setup guidance will be provided. Students using Windows
may need to install additional tools to obtain the command-line environment
used in class.

Readings will include official documentation and short selections on data,
visualization, design, computation, machine learning, and model
documentation. Additional readings will be distributed through the course
site or Canvas.

\section{Course Calendar}

The calendar describes the planned arc of the semester and may change in
response to the pace and interests of the class. Changes will be announced in
class and through the course site or Canvas.

\renewcommand{\arraystretch}{1.18}
\begin{tabularx}{\textwidth}{@{}p{0.12\textwidth}p{0.22\textwidth}X@{}}
\toprule
\textbf{Date} & \textbf{Topic} & \textbf{Questions and Activities} \\
\midrule

Aug.\ 31 &
\textbf{Useful Data} &
What problems are interesting to us? What data could help? Data hunt;
datasets as evidence; what is missing? Introduction to files, the Unix shell,
Git, and GitHub. \\

Sep.\ 7 &
\textbf{No Class} &
Labor Day. \\

Sep.\ 14 &
\textbf{Python as a Tool} &
Jupyter; expressions, variables, strings, numbers, Boolean values, and
computational questions about data. \\

Sep.\ 21 &
\textbf{Programming with Data} &
Collections, conditionals, iteration, functions, decomposition, and
debugging. \\

Sep.\ 28 &
\textbf{Data as Material} &
\texttt{pandas}; Series and DataFrames; reading CSV files; selecting,
filtering, and inspecting data. Project I introduced. \\

Oct.\ 5 &
\textbf{Transforming Data} &
Missing values, types, strings, dates, grouping, aggregation, and the choices
involved in cleaning data. \\

Oct.\ 12 &
\textbf{Seeing Data} &
Visualization; visual encoding; exploratory and explanatory graphics;
asking what a visualization makes visible. \\

Oct.\ 19 &
\textbf{Getting Data} &
HTTP, APIs, JSON, \texttt{requests}, public data sources, and light web data
acquisition. Project I due. \\

Oct.\ 26 &
\textbf{Designing with Data} &
Critique, annotation, narrative, uncertainty, audience, and visual argument.
Project II introduced. \\

Nov.\ 2 &
\textbf{Interactive Data} &
Streamlit fundamentals; widgets, layout, interaction, and designing an
interface around data. \\

Nov.\ 9 &
\textbf{Models and Prediction I} &
Features and targets; training and test data; regression, classification, and
scikit-learn. \\

Nov.\ 16 &
\textbf{Models and Prediction II} &
Evaluation, prediction, overfitting, failure, and model limitations.
Project II due. \\

Nov.\ 23 &
\textbf{Neural Networks and Pretrained Models} &
Training versus inference; neural networks, embeddings, and transformers at a
conceptual level. Final project introduced. \\

Nov.\ 30 &
\textbf{Hugging Face} &
The Model Hub, model cards, pipelines, selecting models, testing inputs and
outputs, and examining limitations. \\

Dec.\ 7 &
\textbf{Designing AI Applications} &
Hugging Face + Streamlit; prototype critique, interface design, failure
modes, testing, and final-project studio. \\

Dec.\ 14 &
\textbf{Final Presentations} &
Presentation and critique of pretrained-model Streamlit applications. \\

\bottomrule
\end{tabularx}

\section{Community Agreements}

Community agreements are classroom expectations for discussion and behavior
that we will create collaboratively. They provide a foundation for an
inclusive learning environment through shared expectations, transparency,
and responsibility. We will establish our agreements early in the semester
and revisit them as needed.

\section{Course Policies}

\subsection{Responsibility}

Students are responsible for assignments and course material even when
absent. If you miss class, consult the course site and communicate with the
instructor or classmates about what you missed. Because much of our work is
cumulative, falling behind on technical material can quickly make later work
more difficult.

\subsection{Participation}

This is a studio-like computational course. Participation includes coming to
class prepared and on time, contributing to discussions and critiques,
working productively with classmates, attempting in-class problems, asking
questions, helping create a productive classroom environment, and engaging
with the process of writing and debugging code.

\subsection{Canvas and Course Materials}

Canvas and the course website will be used for announcements, assignments,
readings, links, and course materials. Students should check for updates
before each class meeting.

\section{Attendance}

At Lang College, regular attendance and participation are important because
learning is both individual and collective. For a course meeting once per
week, the Lang attendance policy calls for a conversation among the student,
instructor, and Student Success Advisor after two absences; more than three
absences normally requires withdrawal or a failing grade, subject to an
agreed attendance plan and applicable university policies.

Absences are counted from the first day a student is enrolled. Students
should communicate with the instructor when circumstances affect their
ability to attend or participate.

\section{Academic Integrity}

Students are responsible for understanding The New School's policies on
academic honesty and integrity and for appropriately distinguishing their own
work from the work of others. This applies to code and computational work as
well as writing, research, presentations, images, and other submitted
materials.

Academic integrity in a programming course includes accurately representing
the sources of code and assistance, following assignment-specific rules about
collaboration and AI-assisted tools, and being able to explain work submitted
as your own.

See The New School University Policies and the University Learning Center for
the current university policy and guidance.

\section{University Resources}

The New School provides resources including University Libraries, the
University Learning Center, Student Disability Services, Archives and Special
Collections, food assistance, health and wellness services, the Student
Ombuds Office, and Financial Aid.

The University Learning Center provides support in areas including software,
computer programming, writing, mathematics, oral presentations, and time
management. Students who need disability-related accommodations should work
with Student Disability Services through the university's established
process.

Current links, contact information, and university policy language are
available through The New School website and Canvas.

\section{A Note on the Calendar}

The topics, project dates, and readings in this syllabus represent our plan
for the semester. Data and programming rarely proceed exactly according to
plan, and neither does a good course. We may adjust the schedule to spend more
time on ideas that require it, respond to student interests, or incorporate
useful datasets and tools that emerge during the semester. Any substantive
changes will be communicated clearly.

\end{document}
